% Generated from chapter-13.md - do not edit.

\chapter{Spirituality}

\section{Turning to God First}

I think that spirituality---my relationship with God---may be one of the most important ``better thoughts'' of all.

One of my best thoughts is simply this:

\emph{I am not alone.}\\
\emph{I can turn to God right now.}

That is what I wish I would do more consistently.

Like many people, I have blown hot and cold in my spiritual life.

I started out fairly devout. I went to Sunday school as a child, and I took faith seriously. One Sunday, a missionary who worked in Africa spoke at our church. I came home convinced that I should go to Africa and help people. It took me years to let go of the guilt I felt when I realized that was not my calling.\index{Africa}

Instead, my gifts turned out to be in teaching, counseling, and helping people communicate---through parent effectiveness training, transactional analysis, couple communication, marriage enrichment, and happiness work. That was where I could truly be useful.\index{Transactional Analysis}\index{Marriage Enrichment}

Then, when our son Jeff died, everything changed.\index{Jeff}

I became very angry at God. I could still talk to Jesus, but I was filled with pain and confusion. Over time, I became a disbeliever and even an atheist. I reasoned that if there were a God, He would never allow my beautiful sixteen-year-old son to die.

It wasn't until Rick became my sweetheart that I slowly turned back toward God and Jesus.

Now I pray, and I believe in prayer. And yet, I still waver sometimes. What feels most solid to me is this: I believe that love is the greatest force in the world, and that goodness is the highest prize to seek.

I believe in love and goodness.

But I find that I can't really pray to ``love and goodness'' alone. Since I see Jesus and God as the living spirit of love and goodness, that is where I direct my prayers.

Today, I want a relationship with God that is sincere and alive---more than formal religion. I want faith without pressure, without pretending, without fear.

Just honesty.

\section{What Is Spirituality?}%
\index{spirituality}\nobreak%

As I've wondered about spirituality, I've talked with friends.

Suzie told me that spirituality is communication with God---contemplating, pondering, and thinking about Jesus. She wants to live with gratitude and humility for everything we've been given: friends, family, guidance, and meaningful work.

Rick says our spiritual lives \emph{are} our lives. He often says,  ``I need to talk to Jesus each day, all day.''

I love that.

It's beautiful. Simple. Honest. Grounded.

He knows where his comfort and strength come from.

Cindy once said she thought the soul lived in the brain. Karo mentioned the movie \emph{21 Grams}, which suggests the soul has physical weight. I learned that this idea came from a small, unreliable experiment over a hundred years ago. Scientists could never confirm it.\index{Cindy}\index{Karo}\index{21 Grams}

It's a lovely metaphor.

Not solid proof.

\section{Where Is the Soul?}%
\index{soul}\nobreak%

I began thinking about music.

Where is music?

In the piano?\\
In the air?\\
In your ear?\\
In your brain?

Yes.\\
And no.

Music happens when everything comes together.

I think the soul is like that.

It is everywhere you are.

When you laugh---it's there.\\
When you pray---it's there.\\
When you cry---it's there.\\
When you listen---it's there.\\
When you forgive---it's there.

The soul isn't stored.

It's expressed.

Across religions and cultures, there is surprising agreement about this.

Most spiritual traditions say the soul is not located in one physical place.

It is presence.\\
Awareness.\\
The breath of God.\\
Spirit.\\
Essence.\\
Light.

In Christianity, we are told, ``The Spirit dwells within you.''

Not in one organ.\\
Not in one corner of the brain.

Within you.

In Judaism, the soul is called ``the breath of God.''

In Buddhism, awareness is understood as experience itself.

In Hinduism, the soul is part of the universal spirit.

Rick's understanding fits beautifully with this. Spirit is not ``somewhere.'' It isn't something you can point to.

It is what happens when body, breath, love, memory, and God's presence come together.

\section{My Better Thought}

For me, spirituality comes down to this:

Turning to God first.

Before worry.\\
Before distraction.\\
Before fear.

When I remember to pray---even simply---I feel steadier and kinder. I remember what matters. I remember who I am.

My better thought is:

\emph{I don't have to do life alone.}\\
\emph{I can talk to God right now.}\\
\emph{I can choose love and goodness today.}

That is enough to guide me.

\section{Drifting Off Course --- and Finding My Way Back}%
\index{way back}\index{drifting}\nobreak%

In my spiritual life, I often feel clear and grounded. I feel close to God. I feel guided by love and goodness. I feel as if I know what matters and how I want to live.

At those times, I think, \emph{Yes. This is it. This is the right direction for me.}

And then . . .  sometimes, I drift.

Not in any dramatic way.\\
Not because I suddenly stop caring.\\
Just quietly, almost without noticing.

I might start watching something that leaves me feeling heavy instead of peaceful.\\
I might eat in a way that doesn't really nourish me.\\
I might distract myself instead of listening to my heart.

Later, I notice: \emph{I don't feel as light. I don't feel as connected. Something feels off.}

For a long time, I wondered why this happened. Why, when I know what feels good and right for me, do I sometimes move away from it?

I used to think it meant I was weak, undisciplined, or failing in some way.

Now I see it differently.

I think of myself as a small boat on the ocean.

Most days, I know where I'm headed.\\
My heart is set on love, kindness, faith, and peace.

But there are winds.\\
There are currents.\\
There are days when I am tired, lonely, worried, or overwhelmed.

And slowly, gently, I drift.

Not because I am bad.\\
Because I am human.

The important thing is not that I never drift.

The important thing is that I notice.

Somewhere inside, I feel a quiet signal:

\emph{This doesn't feel nourishing.}\\
\emph{This isn't bringing me closer to love.}\\
\emph{I miss my good feeling.}

That gentle discomfort is not guilt.\\
It is wisdom.\\
It is God's kindness within me, whispering, \emph{``Come back, sweetheart.''}

When I learned to listen to that voice, everything changed.

Instead of scolding myself, I began asking:

\emph{What am I needing right now?}\\
\emph{Am I tired? Lonely? Overstimulated?}\\
\emph{Do I need rest, comfort, prayer, or connection?}

Often, when I answer those questions honestly, I naturally turn back toward what helps me feel alive again.

A quiet moment.\\
A prayer.\\
A walk.\\
A good conversation.\\
A nourishing meal.\\
A gentle show.\\
A grateful thought.

Coming back does not require punishment.\\
It only requires kindness.

One of my best ``better thoughts'' is this:

\emph{I am allowed to drift sometimes.}\\
\emph{I am always allowed to return.}\\
\emph{Love is still my direction.}

My spiritual life is not about perfection.

It is about relationship.\\
With God.\\
With myself.\\
With life.

And relationships grow through forgiveness, patience, and starting again.

Every time I notice and return, I strengthen my trust in myself.

I learn that I am safe in God's love.\\
I learn that I cannot fall out of grace.\\
I learn that I can always choose a kinder thought.

And that is enough.

\section{About God Being Here}

What if God is not somewhere else?

What if God is not watching from a distance, keeping score, deciding when to intervene?

What if God is the depth of this very moment?

The aliveness in the air.\\
The breath moving in and out.\\
The quiet pulse of Being underneath everything.

What if the sacred is not ``out there'' but woven through every person, every tree, every sorrow, every joy?

And still --- greater than all of it.

Like the ocean is in every wave,\\
but the ocean is more than any one wave.

Like music moving through a piano,\\
but not confined to the piano.

Transcendence would not mean far away.\\
It would mean unfathomably near.

So near that we cannot grasp it.\\
So present that we overlook it.

The better thought is this:\\
Nothing is separate.\\
No moment is empty.\\
No person is without depth.

We are living inside a mystery that holds us.\\
And that mystery is not distant ---\\
it is breathing us.

\textbf{Where do I see God's presence now?}

\section{Looking for Love and Goodness}%
\index{looking for love}\nobreak%

Here is something I discovered almost by accident.

I realized that what I really want in life is very simple.\\
I want to see \textbf{love and goodness}.

I believe love and goodness are the finest things in the world.\\
I also believe they are everywhere.

But there is a problem.

My mind can easily fall into worry. I start planning and thinking about what might go wrong. There is a time for planning, of course. Planning can be wise. But sometimes planning turns into endless speculation, and speculation slowly turns into worry.

When I notice that happening, I remember something that helps me break the spell.

I ask myself a question.

Sometimes I say:

``God, would you show me some love and goodness?''

Or I ask myself:

``Where can I see love and goodness right now?''

When I ask that question, something interesting happens.

If I simply look around my room without thinking about it, my negative bias often goes into operation and I notice something ugly or irritating.

But if I look \textbf{for love and goodness}, I begin to see it.

I might see it in the light coming through the window.\\
Or in a memory of someone who cared about me.\\
Or in the thought that I could offer kindness to someone today.

Sometimes I even ask another question:

``How could I give someone love and goodness right now?''

When I ask that, ideas begin to show up.

\textbf{A better thought:}\\
Love and goodness may be everywhere, but I see them more easily when I go looking for them.

\textbf{A question for you:}\\
If you looked around right now with the intention of seeing love and goodness, what might you notice?
